%%% FIELD prop-statement
Suppose \(h_n\le 1/32\), with \(a_\star=1-v\) and \(b=1-v/2\).  Then \(J_n\ge2\), the cutoff locations in \(\mathcal W_n\) lie in \([1/4,3/4]\), and the supports of distinct compliance profiles are disjoint up to Lebesgue-null endpoints.  The common potential-outcome table satisfies
\[
q_{00}(x)=q_{11}(x)=\frac v4,
\qquad
\frac v4\le q_{01}(x),q_{10}(x)\le1-\frac{3v}{4},
\qquad
\sum_{u_0,u_1\in\{0,1\}}q_{u_0,u_1}(x)=1.
\]
All twelve latent masses defining every \(P_{j,\ell}\in\mathcal W_n\) are nonnegative and sum to one; the always-taker and never-taker masses are strictly positive everywhere, and the complier masses are strictly positive wherever \(w_{j,\ell}>0\).  No H\"older restriction on these latent cells is used.  The observable one-sided extensions are
\[
\mu_{D,P_{j,\ell}}^{-}(x)=\frac{1-w_{j,\ell}(x)}2,
\qquad
\mu_{D,P_{j,\ell}}^{+}(x)=\frac{1+w_{j,\ell}(x)}2,
\]
\[
\mu_{Y,P_{j,\ell}}^{-}(x)=\frac12-\frac{g_\star(x)w_{j,\ell}(x)}2,
\qquad
\mu_{Y,P_{j,\ell}}^{+}(x)=\frac12+\frac{g_\star(x)w_{j,\ell}(x)}2.
\]
They belong to the radius-\(M\), order-\(s\) H\"older ball.  The potential reduced forms are continuous at \(c_{j,\ell}\), and
\[
d_{P_{j,\ell}}=\kappa_n,
\qquad
\theta_{P_{j,+}}=a_\star=1-v,
\qquad
\theta_{P_{j,-}}=-a_\star=-(1-v).
\]
Writing \(p_{j,\ell}^{D,Y}\) and \(p_{j,\ell}^{D}\) for the conditional observable cell probabilities, one has
\[
m_{d,y}(x)\ge\frac v2,
\qquad
p_{j,\ell}^{D,Y}(x,d,y)
=p_{j,\ell}^{D}(x,d)m_{d,y}(x)
\ge\frac{7v}{32}>0.
\]
On either observable side,
\[
\operatorname{Var}_{P_{j,\ell}}(D\mid X=x)=\frac{1-w_{j,\ell}(x)^2}{4}>0,
\]
\[
\operatorname{Var}_{P_{j,\ell}}(Y\mid X=x)
=\frac{1-g_\star(x)^2}{4}
+g_\star(x)^2\operatorname{Var}_{P_{j,\ell}}(D\mid X=x)
=\frac{1-g_\star(x)^2w_{j,\ell}(x)^2}{4}
\ge\frac{63}{256}>v,
\]
and
\[
\left|\operatorname{corr}_{P_{j,\ell}}(Y,D\mid X=x)\right|
=\frac{|g_\star(x)|\sqrt{1-w_{j,\ell}(x)^2}}
       {\sqrt{1-g_\star(x)^2w_{j,\ell}(x)^2}}
\le |g_\star(x)|\le1-v.
\]
The formula remains valid when \(w_{j,\ell}(x)=0\); treatment variance is never degenerate in this witness.  Consequently \(\mathcal W_n\subset\mathcal P_n\).  Moreover, with \(R_\star\) as the outcome-invariant reference,
\[
\frac{dP_{j,\ell}}{dR_\star}
=1+(2D-1)\operatorname{sgn}(X-c_{j,\ell})w_{j,\ell}(X)
\quad R_\star\text{-almost surely},
\]
so every product likelihood ratio is measurable with respect to \(\sigma((X_i,D_i)_{i=1}^n)\).  Finally,
\[
\int_{-1}^{1}\psi(u)^2\,du=\frac{256}{315},
\qquad
E_{R_\star}\!\left[
\frac{dP_{j,\ell}}{dR_\star}
\frac{dP_{k,r}}{dR_\star}
\right]=1
\]
for distinct \((j,\ell)\ne(k,r)\), and, for
\(Q_{\ell,n}=J_n^{-1}\sum_{j=1}^{J_n}P_{j,\ell}^{\otimes n}\),
\[
\chi^2(Q_{\ell,n},R_\star^{\otimes n})
=\frac{\left(1+\frac{256}{315}\kappa_n^2h_n\right)^n-1}{J_n}.
\]
%%% FIELD prop-proof
Write \(a=a_\star=1-v\) and \(b=1-v/2\).  We verify each model property and likelihood calculation used by the lower bound.

First, \(h_n\le1/32\) implies \((16h_n)^{-1}\ge2\), and hence
\[
J_n=\left\lfloor(16h_n)^{-1}\right\rfloor\ge2.
\tag{1}
\]
For \(1\le j\le J_n\), the location formulas in \(\mathrm{def:outcome\mbox{-}invariant\mbox{-}family}\) give
\[
c_{j,+}-h_n=\frac14+2(j-1)h_n\ge\frac14,
\qquad
c_{j,+}+h_n=\frac14+2jh_n\le\frac38,
\tag{2}
\]
and similarly
\[
[c_{j,-}-h_n,c_{j,-}+h_n]\subset[5/8,3/4].
\tag{3}
\]
Centers within either label are separated by \(2h_n\), and the two labels use the disjoint plateau intervals in (2)--(3).  Because \(\psi\) is supported on \([-1,1]\), distinct supports meet at most at endpoints.  Uniformity of \(X\) makes those endpoints null, and therefore
\[
w_{j,\ell}(x)w_{k,r}(x)=0
\quad\text{for Lebesgue-almost every }x
\quad\text{whenever }(j,\ell)\ne(k,r).
\tag{4}
\]

Next, \(B'(t)=-30t^2(1-t)^2\le0\), \(B(0)=1\), and \(B(1)=0\).  Thus \(\mathrm{def:common\mbox{-}outcome\mbox{-}profile}\) gives \(|g_\star(x)|\le a=1-v\).  By the corrected construction in \(\mathrm{def:common\mbox{-}outcome\mbox{-}table}\),
\[
q_{00}=q_{11}=\frac v4,
\qquad
q_{01}=\frac{1-v/2+g_\star}{2},
\qquad
q_{10}=\frac{1-v/2-g_\star}{2}.
\tag{5}
\]
The four entries sum to one.  Since \(|g_\star|\le1-v\),
\[
\frac v4
=\frac{1-v/2-(1-v)}2
\le q_{01},q_{10}
\le\frac{1-v/2+(1-v)}2
=1-\frac{3v}{4}.
\tag{6}
\]
Together with \(v>0\), (5)--(6) show that every common-table cell is strictly positive.

The quartic profile obeys \(0\le\psi\le1\).  Hence \(0\le w_{j,\ell}\le\kappa_n\le1/8\), and the latent masses
\[
\pi_{C,u_0,u_1,P_{j,\ell}}=w_{j,\ell}q_{u_0,u_1},
\qquad
\pi_{A,u_0,u_1,P_{j,\ell}}
=\pi_{N,u_0,u_1,P_{j,\ell}}
=\frac{1-w_{j,\ell}}2q_{u_0,u_1}
\tag{7}
\]
are nonnegative.  The always-taker and never-taker cells are bounded below by
\[
\frac{1-1/8}{2}\frac v4=\frac{7v}{64}>0,
\tag{8}
\]
and a complier cell is positive wherever \(w_{j,\ell}>0\).  Summing (7) over response types and the four potential-outcome cells gives
\[
\sum_{u_0,u_1}q_{u_0,u_1}
\left(w_{j,\ell}+\frac{1-w_{j,\ell}}2+
                    \frac{1-w_{j,\ell}}2\right)=1.
\tag{9}
\]
This proves validity of the latent law using positivity alone; no smoothness of an individual latent cell is imposed or used.

We next verify the observable and potential reduced forms defining \(\mathcal P_n\).  The response-type weights in (7) imply
\[
E[D(z)\mid X=x]=\frac{1-w_{j,\ell}(x)}2+z\,w_{j,\ell}(x).
\tag{10}
\]
Because the potential-outcome table is common to every response type, (5) gives
\[
E[Y(0)\mid X=x]=\frac{1-g_\star(x)}2,
\qquad
E[Y(1)\mid X=x]=\frac{1+g_\star(x)}2.
\tag{11}
\]
Combining (10)--(11) yields
\[
E[Y(D(z))\mid X=x]
=\frac12+g_\star(x)\left(z-\frac12\right)w_{j,\ell}(x).
\tag{12}
\]
Equations (10)--(12), consistency, and \(Z=\mathbf 1\{X\ge c_{j,\ell}\}\) give the four displayed observable extensions in the statement.

It remains to check their H\"older radii.  On \((-1,1)\),
\[
\psi'(u)=-4u(1-u^2),
\qquad
\psi''(u)=-4+12u^2.
\tag{13}
\]
Both \(\psi\) and \(\psi'\) join continuously to zero at \(\{-1,1\}\).  Consequently
\[
\lVert\psi'\rVert_\infty\le2,
\qquad
\operatorname{Lip}(\psi')\le8.
\tag{14}
\]
When \(s=1\), the identity \(h_n=32\kappa_n/M\) and (14) give
\[
\operatorname{Lip}(w_{j,\ell})
\le\frac{2\kappa_n}{h_n}=\frac{M}{16}.
\tag{15}
\]
When \(1<s\le2\), put \(\beta=s-1\).  The bounded, globally Lipschitz function \(\psi'\) has \(\beta\)-H\"older seminorm at most \(8\): for arguments within unit distance this follows from (14), and for arguments farther apart it follows from \(2\lVert\psi'\rVert_\infty\le4\).  Since \(h_n^s=32\kappa_n/M\) and \(h_n\le1\), scaling gives
\[
\lVert w_{j,\ell}'\rVert_\infty
\le\frac{2\kappa_n}{h_n}
=\frac{M}{16}h_n^{s-1}\le\frac{M}{16},
\qquad
[w_{j,\ell}']_{s-1}
\le\frac{8\kappa_n}{h_n^s}=\frac{M}{4}.
\tag{16}
\]
By (2)--(3), \(g_\star=a\) on the support of every \(+\) profile and \(g_\star=-a\) on the support of every \(-\) profile.  Since each profile vanishes off its support,
\[
g_\star w_{j,+}=a w_{j,+},
\qquad
g_\star w_{j,-}=-a w_{j,-}
\quad\text{on }[0,1].
\tag{17}
\]
Thus, for \(s=1\), the treatment and outcome extensions have Lipschitz constants at most \(M/32\) and \(aM/32\), respectively.  For \(1<s\le2\), their derivative suprema are bounded by the same quantities and their derivative H\"older seminorms are at most \(M/8\) and \(aM/8\), respectively.  Since \(a<1\), all four extensions have supremum norm at most \(1/2+\kappa_n/2\le9/16\).  Even under the convention that the radius bounds the sum of these norm components,
\[
\frac9{16}+\frac{M}{32}+\frac{M}{8}\le M
\tag{18}
\]
because \(M\ge2\); deleting the derivative-seminorm term when \(s=1\) only strengthens the inequality.  Hence the extensions satisfy \(\mathrm{ass:observable\mbox{-}holder\mbox{-}extensions}\).  Equations (10), (12), and (17) are continuous at \(c_{j,\ell}\), so \(\mathrm{ass:potential\mbox{-}reduced\mbox{-}form\mbox{-}continuity}\) also holds.

Since \(w_{j,\ell}(c_{j,\ell})=\kappa_n>0\), the observable treatment extensions give
\[
d_{P_{j,\ell}}=\kappa_n.
\tag{19}
\]
The outcome extensions and (17) give
\[
y_{P_{j,+}}=a\kappa_n,
\qquad
y_{P_{j,-}}=-a\kappa_n.
\tag{20}
\]
Division by the positive denominator in (19) is therefore legitimate, and the definition of \(\theta_P\) yields the causal labels \(1-v\) and \(-(1-v)\).  The cutoff is interior by (2)--(3), \(X\) is uniform by construction, and taking independent copies gives the required iid sample.

For nondegeneracy and outcome invariance, the common table makes the conditional law of \(Y\) given \((D=d,X=x)\) Bernoulli with success probability
\[
E[Y\mid D=d,X=x]=\frac12+g_\star(x)\left(d-\frac12\right).
\tag{21}
\]
Every success or failure probability in (21) is at least
\[
\frac{1-|g_\star(x)|}{2}\ge\frac v2.
\tag{22}
\]
On either side of the cutoff, the two conditional treatment probabilities are \((1+w_{j,\ell})/2\) and \((1-w_{j,\ell})/2\), in an order determined by the side, so each is at least \(7/16\).  Therefore the common-table factorization gives
\[
p_{j,\ell}^{D,Y}(x,d,y)
=p_{j,\ell}^{D}(x,d)m_{d,y}(x)
\ge\frac{7v}{32}>0.
\tag{23}
\]
This strict positivity permits density cancellation below; no comparison of the joint-cell floor with the model constant \(v\) is needed or asserted.

Let \(w=w_{j,\ell}(x)\) and \(g=g_\star(x)\).  On either side, \(D\mid X=x\) is Bernoulli with parameter \((1+w)/2\) or \((1-w)/2\), and hence
\[
V_D:=\operatorname{Var}(D\mid X=x)=\frac{1-w^2}{4}\ge\frac{63}{256}>0.
\tag{24}
\]
The conditional variance identity and (21) give
\[
\operatorname{Var}(Y\mid X=x)
=\frac{1-g^2}{4}+g^2V_D
=\frac{1-g^2w^2}{4}
\ge\frac{63}{256}>\frac1{64}>v,
\tag{25}
\]
where the first lower bound uses \(|g|\le1\) and \(w\le1/8\), and the last uses \(0<v<1/64\).  Likewise,
\[
\operatorname{Cov}(Y,D\mid X=x)=gV_D.
\tag{26}
\]
Because (24) is strictly positive, (24)--(26) imply the exact formula
\[
|\operatorname{corr}(Y,D\mid X=x)|
=\frac{|g|\sqrt{1-w^2}}{\sqrt{1-g^2w^2}}
\le |g|\le1-v.
\tag{27}
\]
Indeed, \(g^2\le1\) gives \(1-w^2\le1-g^2w^2\).  Formula (27) also applies at \(w=0\), where its square-root ratio equals one.  No treatment-degenerate point occurs in this witness by (24); if such a point occurred in a different law, (26) would give zero covariance and the convention in \(\mathrm{ass:observable\mbox{-}correlation\mbox{-}separation}\) would apply without forming a correlation coefficient.  Equations (2)--(27) verify every member property in \(\mathrm{def:model\mbox{-}class}\), so
\[
\mathcal W_n\subset\mathcal P_n.
\tag{28}
\]

Finally, under \(R_\star\), \(D\) is conditionally fair given \(X\), while (21) is the same strictly positive conditional outcome margin as under every \(P_{j,\ell}\).  Cancelling that margin in the full observable density ratio gives, outside the null event \(X=c_{j,\ell}\),
\[
L_{j,\ell}:=\frac{dP_{j,\ell}}{dR_\star}
=1+(2D-1)\operatorname{sgn}(X-c_{j,\ell})w_{j,\ell}(X).
\tag{29}
\]
Thus the one-unit and product likelihood ratios are measurable with respect to the corresponding \((X,D)\) sigma-fields.  Direct integration gives
\[
\int_{-1}^{1}\psi(u)^2\,du
=2\left(1-\frac43+\frac65-\frac47+\frac19\right)
=\frac{256}{315}.
\tag{30}
\]
Conditioning on \(X\) under \(R_\star\) in (29), and using \(E_{R_\star}[2D-1\mid X]=0\), yields
\[
E_{R_\star}[L_{j,\ell}L_{k,r}]
=1+\int_0^1
\operatorname{sgn}(x-c_{j,\ell})
\operatorname{sgn}(x-c_{k,r})
w_{j,\ell}(x)w_{k,r}(x)\,dx.
\tag{31}
\]
For distinct indices, (4) makes (31) equal to one.  For equal indices, the substitution \(u=(x-c_{j,\ell})/h_n\) and (30) give
\[
E_{R_\star}[L_{j,\ell}^2]
=1+\frac{256}{315}\kappa_n^2h_n.
\tag{32}
\]
Independence raises (31)--(32) to the \(n\)-th power.  Expanding the square of the equally weighted mixture likelihood ratio gives \(J_n\) diagonal terms and \(J_n(J_n-1)\) off-diagonal terms.  Subtracting one therefore yields
\[
\chi^2(Q_{\ell,n},R_\star^{\otimes n})
=\frac{\left(1+\frac{256}{315}\kappa_n^2h_n\right)^n-1}{J_n},
\]
which completes the proof.
%%% FIELD theorem-statement
With \(a_\star=1-v\), put
\[
C_\psi=\frac{256}{315},
\qquad
C_h=\left(\frac{32}{M}\right)^{1/s},
\qquad
\lambda_w=(sC_\psi C_h)^{-1}
=\frac{315}{256s}\left(\frac{M}{32}\right)^{1/s}.
\]
If \(h_n\le1/32\), then the high-separation outcome-invariant witness gives the finite-sample bound
\[
L_n(\kappa_n)
\ge2(1-v)\left[
1-2\alpha-
\sqrt{\frac{\left(1+C_\psi\kappa_n^2h_n\right)^n-1}{J_n}}
\right]_+.
\]
Along every admissible triangular sequence,
\[
\frac1s\log(1/\kappa_n)-C_\psi C_hT_n\longrightarrow+\infty
\quad\Longrightarrow\quad
\liminf_{n\to\infty}L_n(\kappa_n)
\ge2(1-v)(1-2\alpha).
\]
In particular,
\[
\limsup_{n\to\infty}
\frac{T_n}{\log(1/\kappa_n)}<\lambda_w
\quad\Longrightarrow\quad
\liminf_{n\to\infty}L_n(\kappa_n)
\ge2(1-v)(1-2\alpha).
\]
For \(s=1\), \(\lambda_w=315M/8192\).  The sufficient regime contains sequences on which the known-cutoff minimax length shrinks to zero: for every \(0<\gamma<(2+1/s)^{-1}\), the eventual choice \(\kappa_n=\rho_n(\log n)^\gamma\) satisfies the first displayed limit and \(L_n^{\mathrm{known}}(\kappa_n)\to0\).  The constant \(\lambda_w\) is an explicit one-sided witness threshold; this theorem does not claim that it is the exact critical frontier.
%%% FIELD theorem-proof
Let \(I_n\) be an arbitrary measurable connected interval rule satisfying the uniform coverage condition in \(\mathrm{def:minimax\mbox{-}length}\).  By \(\mathrm{prop:outcome\mbox{-}invariant\mbox{-}witness\mbox{-}high\mbox{-}separation}\), each component of \(Q_{+,n}\) is in \(\mathcal P_n\) and has causal target \(a_\star=1-v\), while each component of \(Q_{-,n}\) is in \(\mathcal P_n\) and has target \(-a_\star=-(1-v)\).  Averaging the componentwise honesty inequalities gives
\[
Q_{+,n}\{a_\star\in I_n\}\ge1-\alpha,
\qquad
Q_{-,n}\{-a_\star\in I_n\}\ge1-\alpha.
\tag{33}
\]
For probability measures \(Q\ll R\), Cauchy--Schwarz applied to the density ratio proves
\[
\operatorname{TV}(Q,R)
=\frac12E_R\left|\frac{dQ}{dR}-1\right|
\le\frac12\sqrt{E_R\left(\frac{dQ}{dR}-1\right)^2}
=\frac12\sqrt{\chi^2(Q,R)}.
\tag{34}
\]
The triangle inequality for total variation, (34), and the equal chi-square identities for the two labels yield
\[
\operatorname{TV}(Q_{+,n},Q_{-,n})
\le\operatorname{TV}(Q_{+,n},R_\star^{\otimes n})
   +\operatorname{TV}(Q_{-,n},R_\star^{\otimes n})
\le
\sqrt{\frac{\left(1+C_\psi\kappa_n^2h_n\right)^n-1}{J_n}}.
\tag{35}
\]
Transfer the second event in (33) from \(Q_{-,n}\) to \(Q_{+,n}\) by the definition of total variation, and intersect it with the first event in (33).  The union bound and (35) give
\[
Q_{+,n}\{-a_\star,a_\star\in I_n\}
\ge1-2\alpha-
\sqrt{\frac{\left(1+C_\psi\kappa_n^2h_n\right)^n-1}{J_n}}.
\tag{36}
\]
A connected subset of \(\Theta=[-1,1]\) containing both labels has Lebesgue length at least \(2a_\star=2(1-v)\).  Since \(Q_{+,n}\) is the uniform mixture of in-class product laws,
\[
\sup_{P\in\mathcal P_n}E_P[|I_n|]
\ge E_{Q_{+,n}}[|I_n|]
\ge2(1-v)\left[
1-2\alpha-
\sqrt{\frac{\left(1+C_\psi\kappa_n^2h_n\right)^n-1}{J_n}}
\right]_+.
\tag{37}
\]
Taking the infimum over all honest rules proves the finite-sample assertion.

We now evaluate (37).  The definitions of \(h_n\), \(C_h\), and \(T_n\) imply
\[
h_n=C_h\kappa_n^{1/s},
\qquad
n\kappa_n^2h_n=C_hT_n.
\tag{38}
\]
Whenever \(h_n\le1/32\), the number \(x=(16h_n)^{-1}\) is at least two.  Since \(\lfloor x\rfloor\ge x/2\) for \(x\ge2\),
\[
J_n\ge\frac1{32h_n}.
\tag{39}
\]
Using \(1+x\le e^x\), equations (38)--(39) bound the chi-square term in (37) by
\[
\frac{\left(1+C_\psi\kappa_n^2h_n\right)^n-1}{J_n}
\le32h_n\exp\{C_\psi n\kappa_n^2h_n\}
=32C_h\kappa_n^{1/s}\exp\{C_\psi C_hT_n\}.
\tag{40}
\]
The logarithm of the final expression is
\[
\log(32C_h)-\frac1s\log(1/\kappa_n)+C_\psi C_hT_n.
\tag{41}
\]
Suppose
\[
\frac1s\log(1/\kappa_n)-C_\psi C_hT_n\longrightarrow+\infty.
\tag{42}
\]
Because \(T_n>0\), condition (42) implies \(\log(1/\kappa_n)\to\infty\).  Hence \(\kappa_n\to0\), so \(h_n\to0\) and (39)--(41) apply eventually.  Equation (41) tends to \(-\infty\), and thus the square-root term in (37) tends to zero.  We conclude that
\[
\liminf_{n\to\infty}L_n(\kappa_n)
\ge2(1-v)(1-2\alpha).
\tag{43}
\]

For the ratio corollary, suppose
\[
\limsup_n\frac{T_n}{\log(1/\kappa_n)}<\lambda_w.
\tag{44}
\]
Condition (44) forces \(\kappa_n\to0\).  Otherwise, along a subsequence bounded away from zero, \(T_n=n\kappa_n^{2+1/s}\) would diverge at least linearly while \(\log(1/\kappa_n)\) remained bounded, contradicting (44).  Choose \(\varepsilon>0\) so that eventually
\[
\frac{T_n}{\log(1/\kappa_n)}\le\lambda_w-\varepsilon.
\]
Since \(C_\psi C_h\lambda_w=1/s\),
\[
\frac1s\log(1/\kappa_n)-C_\psi C_hT_n
\ge C_\psi C_h\varepsilon\log(1/\kappa_n)
\longrightarrow+\infty.
\tag{45}
\]
Thus (44) implies (42), and (43) proves the stated corollary.  Substitution of \(s=1\) gives
\[
\lambda_w=\frac{315}{256}\frac{M}{32}=\frac{315M}{8192}.
\]

Finally, fix \(0<\gamma<(2+1/s)^{-1}\) and, for all sufficiently large \(n\), set
\[
\kappa_n=\rho_n(\log n)^\gamma
=n^{-s/(2s+1)}(\log n)^\gamma.
\tag{46}
\]
This sequence is positive and eventually at most \(1/8\).  Direct calculation gives
\[
T_n=n\kappa_n^{2+1/s}
=(\log n)^{\gamma(2+1/s)}=o(\log n),
\tag{47}
\]
while
\[
\log(1/\kappa_n)
=\frac{s}{2s+1}\log n-\gamma\log\log n
\asymp\log n.
\tag{48}
\]
Equations (47)--(48) imply (42).  Moreover,
\[
\frac{\rho_n}{\kappa_n}=(\log n)^{-\gamma}\longrightarrow0.
\]
The upper inequality in \(\mathrm{thm:known\mbox{-}cutoff\mbox{-}benchmark}\) therefore gives \(L_n^{\mathrm{known}}(\kappa_n)\to0\), whereas (43) remains in force.  This proves the oracle-shrinking comparison.  No converse or attaining construction is supplied at or above \(\lambda_w\), so the conclusion remains the asserted one-sided witness threshold.
%%% FIELD common-table-construction
Set \(b=1-v/2\).  For \(x\in[0,1]\), define \(q_{00}(x)=q_{11}(x)=(1-b)/2\), \(q_{01}(x)=\{b+g_\star(x)\}/2\), and \(q_{10}(x)=\{b-g_\star(x)\}/2\).
%%% FIELD critical-statement
Using \(\mathsf H_{\mathrm{crit}}\), determine a nonvariational
fixed-\(\alpha\) expression for \(L_n(\kappa_n)\) with a uniform
remainder throughout \(T_n\asymp\log(1/\kappa_n)\), and construct an
honest connected interval attaining that expression.  In particular,
for every \(\lambda\in(0,\infty)\), derive and match the limiting
expected length along sequences satisfying
\(T_n/\log(1/\kappa_n)\to\lambda\).  The result must retain
effect-label testing, agree with the exact \(s=1\) common-outcome
plateau experiment, have a \(\lambda\downarrow0\) converse that
recovers Theorem 2's nonshrinking lower bound, with no matched
below-boundary value claimed without a separate attaining upper bound,
and agree in order with Theorems 1 and 3 as
\(\lambda\uparrow\infty\).  It must strictly extend the known-cutoff
honesty result of Noack and Rothe (2024) by evaluating connected
expected length uniformly over the unknown cutoff \(c_P\).
%%% FIELD critical-partial
The explicit witness already evaluates one side of the critical experiment.  Put \(C_h=(32/M)^{1/s}\) and \(C_\psi=256/315\).  Then \(h_n=C_h\kappa_n^{1/s}\), \(\log J_n=s^{-1}\log(1/\kappa_n)+O(1)\), and
\[
nC_\psi\kappa_n^2h_n=C_\psi C_hT_n.
\]
Thus, along \(T_n/\log(1/\kappa_n)\to\lambda\), the displayed mixture chi-square divergence tends to zero whenever \(C_\psi C_h\lambda<1/s\).  In that subregion the finite-sample argument gives \(\liminf_nL_n(\kappa_n)\ge2(1-v)(1-2\alpha)\).  For \(s=1\), the witness threshold is \(\lambda=315M/8192\).  At equality, the first-order coordinate \(\lambda\) does not determine this mixture's limit: the second-order value of \(C_\psi C_hT_n-s^{-1}\log(1/\kappa_n)\) enters.  Above the witness threshold its chi-square divergence grows and the construction supplies no converse.  No honest connected interval with a matched fixed-\(\alpha\) expected-length constant has been constructed in the remaining critical region.
